%------------------------------------------------------------------------------%
% SPECTRAL THEORY OF PARTIAL DIFFERENTIAL OPERATORS                            %
%------------------------------------------------------------------------------%
% File  : Spectral_Theory_of_Partial_Differential_Operators.tex                %
% Author: Klaus Hoeflinger, Beethovenstrasse 11, D-73117 Wangen                %
% Email : khoeflinger@t-online.de                                              %
%------------------------------------------------------------------------------%

%------------------------------------------------------------------------------%
%                       DOCUMENT CLASS and INCLUDE FILES                       %
%------------------------------------------------------------------------------%
\documentclass[11pt,twoside]{article}
\usepackage{geometry}
\geometry{paperheight=241mm,
          paperwidth=152mm,
          top=22mm,
          bottom=17mm,
          left=20mm,
          right=20mm,
          textwidth=112mm,
          textheight=202mm,
          headheight=10mm,
          headsep=3mm,
          footskip=9mm,
          includehead,
          includefoot,
          marginparsep=0mm,
          marginparwidth=0mm}

\renewcommand{\baselinestretch}{0.945}\normalsize

\AtBeginDocument{\setlength\abovedisplayskip{2mm}}
\AtBeginDocument{\setlength\belowdisplayskip{2mm}}

%------------------------------------------------------------------------------%
% define title formats                                                         %
%------------------------------------------------------------------------------%
\usepackage{titlesec}

\renewcommand{\thesection}{\arabic{section}}
\renewcommand{\sfdefault}{FiraSans}
\renewcommand{\ttdefault}{pcr}
\renewcommand{\rmdefault}{antiqua}

\titleformat{\part}[display]
{\normalfont \large \rmfamily \bfseries \centering}
{\small{\thepart.}}{2pt}{}

\titleformat{\section}[runin]
{\normalfont \rmfamily \bfseries}
{\thesection}{1pt}{.\;}[.]

%\titleformat{\section}
%{\normalfont \bfseries \small \filcenter}
%{\thesection.\;}{0mm}{}

\titleformat{\subsection}[runin]
{\normalfont \itshape}
{}{0mm}{}

\titleformat{\subsubsection}[runin]
{\normalfont \itshape}
{}{}{}{}

\titleformat{\paragraph}[runin]
{\normalfont}
{}{}{}{}

\titleformat{\subparagraph}[runin]
{\normalfont}
{}{}{}{}

\titlespacing*{\part}          { 0pt}{ 0pt}{ 8pt}
\titlespacing*{\section}       { 0pt}{ 7pt}{ 4pt}
\titlespacing*{\subsection}    { 0pt}{14pt}{ 6pt}
\titlespacing*{\subsubsection} { 0pt}{ 6pt}{12pt}
\titlespacing*{\paragraph}     { 0pt}{ 6pt}{12pt}
\titlespacing*{\subparagraph}  { 0pt}{ 0pt}{12pt}

%------------------------------------------------------------------------------%
% define enumerate parameters                                                  %
%------------------------------------------------------------------------------%
\usepackage{enumitem}
\setlength{\parindent}{3pt}
\setlength{\parskip}  {0pt}

%\setlist{noitemsep}

\setlist[enumerate]{
         leftmargin=12pt,
         rightmargin=0pt,
         itemsep=1pt,
         itemindent=3pt,
         topsep=3pt,
         labelsep=4pt,
         partopsep=0pt,
         labelwidth=0pt,
         listparindent=0pt,
         parsep=0pt}

\setlist[itemize]{
         leftmargin=12pt,
         rightmargin=0pt,
         itemsep=0pt,
         itemindent=1pt,
         topsep=1pt,
         labelsep=4pt,
         partopsep=1pt,
         labelwidth=0pt,
         listparindent=0pt,
         parsep=0pt}

%------------------------------------------------------------------------------%
% define packages                                                              %
%------------------------------------------------------------------------------%
\usepackage{upref}
\usepackage{amssymb}
\usepackage{slantsc}
\usepackage{amsmath}
\usepackage{amsthm}
\usepackage{thmtools}
\usepackage{amscd}
\usepackage{verbatim}
\usepackage{latexsym}
\usepackage{multicol}
\usepackage{graphicx}
\usepackage{setspace}
\usepackage[ngerman,english]{babel}
\usepackage[protrusion=true,expansion=true]{microtype}
\usepackage{tikz}
\usetikzlibrary{arrows,arrows.meta,decorations.markings}
\usepackage{mathrsfs}
\usepackage{amsfonts}
\usepackage{datetime2}
\usepackage{textgreek}
\usepackage{hyperref}
\usepackage{pifont}
\usepackage{relsize}
%\usepackage{biblatex}
\relsize{-0.05}
\usepackage[skip=0mm]{parskip}
\usetikzlibrary{decorations.pathmorphing}

\usepackage{mlmodern}
\usepackage[T5,T1]{fontenc}
\usepackage{ragged2e}

%\usepackage{mathabx}
%\usepackage{times}

%\usepackage{fontspec}
%\setmainfont{Nimbus Roman No9 L}
%\usepackage[T1]{fontenc}

%------------------------------------------------------------------------------%
% define page layout                                                           %
%------------------------------------------------------------------------------%
\usepackage{fancyhdr}
\pagestyle{fancy}
\setlength{\headheight}{30.0pt}
\renewcommand{\headrulewidth}{0pt}

\AtBeginDocument{
  \addtolength\abovedisplayskip{-0.0\baselineskip}
  \addtolength\belowdisplayskip{-0.0\baselineskip}
}

\fancyhead{}
\fancyfoot{}
\addtolength{\headwidth}{0.02\marginparwidth}

\makeatletter
\let\ps@plain\ps@fancy
\makeatother

\renewcommand\sectionmark[1]
{
 \def\sectionname{#1}
 \markright{\thesection #1}
}

\makeatletter
\@addtoreset{section}{part}
\makeatother

%------------------------------------------------------------------------------%
% define footnote without number                                               %
%------------------------------------------------------------------------------%
\newcommand{\myfootnote}[1]{%
\renewcommand{\thefootnote}{}%
\footnotetext{#1}%
\renewcommand{\thefootnote}{\arabic{footnote}}%
}

%------------------------------------------------------------------------------%
% define numbering in equations                                                %
%------------------------------------------------------------------------------%
\numberwithin{equation}{section}

%------------------------------------------------------------------------------%
% define newtheorem                                                            %
%------------------------------------------------------------------------------%
\declaretheoremstyle[
headfont=\scshape, 
headindent = 0mm, %\parindent,
%postheadhook = {\hspace{0mm}\newline},
%postheadspace = \newline, 
spaceabove = 3mm, 
spacebelow = 3mm,
numberwithin=section,
name=Theorem]{khMATH}
\declaretheorem[style=khMATH]{theorem}

\declaretheoremstyle[
headfont=\bfseries, 
headindent = 0mm, %\parindent,
%postheadhook = {\hspace{0mm}\newline},
%postheadspace = \newline, 
spaceabove = 3mm, 
spacebelow = 3mm,
numberwithin=part,
name=Theorem]{khMATH1}
\declaretheorem[style=khMATH1]{theorem1}

\declaretheoremstyle[
headfont=\scshape, 
headindent = 0mm, %\parindent,
%postheadhook = {\hspace{0mm}\newline},
%postheadspace = \newline, 
spaceabove = 3mm, 
spacebelow = 3mm,
numberwithin=section,
name=Corollary]{khMATH1}
\declaretheorem[style=khMATH1]{corollary}

%            name         counter     label              numeration-mode
\theoremstyle{plain}
\newtheorem*{introduction}            {\sc Introduction}
\newtheorem {standard}                {}                 [section]
\newtheorem*{definition}              {\sc Definition}
\newtheorem{satz1}                    {\sc Satz}
\newtheorem{lemma}                    {\sc Lemma}
\newtheorem*{proposition}             {\sc Proposition}
%\newtheorem{corollary}               {\sc Corollary}
\newtheorem*{corollar}                {\sc Korollar}
\newtheorem*{remark}                  {\sc Remark}
\newtheorem*{remarks}                 {\sc Remarks}
\newtheorem*{example}                 {Example}
\newtheorem*{examples}                {Examples}
\newtheorem*{theo}                    {\sc Theorem}

%------------------------------------------------------------------------------%
% define tabular                                                               %
%------------------------------------------------------------------------------%
\usepackage{tabularx}
\newcolumntype{L}[1]{>{\raggedright\arraybackslash}p{#1}}
\newcolumntype{C}[1]{>{\centering\arraybackslash}  p{#1}} 
\newcolumntype{R}[1]{>{\raggedleft\arraybackslash} p{#1}} 

%------------------------------------------------------------------------------%
% define \sh, \zB                                                              %
%------------------------------------------------------------------------------%
\def\dsh{{d.\,h.}}
\def\ie{{i.\,e.}}
\def\zB{z.\,B.}
\renewcommand{\abstractname}{ }

%------------------------------------------------------------------------------%
% change default spacing for cases                                             %
%------------------------------------------------------------------------------%
\makeatletter
\renewcommand*\env@cases[1][1.6]{%
  \let\@ifnextchar\new@ifnextchar
  \left\lbrace
  \def\arraystretch{#1}%
  \array{@{}l@{\quad}l@{}}%
}
\makeatother

\newenvironment{rcases}
  {\left.\begin{aligned}}
  {\end{aligned}\right\rbrace}

%------------------------------------------------------------------------------%
% change default for \predisplaypenalty                                        %
%------------------------------------------------------------------------------%
\predisplaypenalty=990
\allowdisplaybreaks \numberwithin{equation}{section}

%------------------------------------------------------------------------------%
% general notations                                                            %
%------------------------------------------------------------------------------%
\def\*{\star}
\def\={\,=\,}
\def\+{\,+\,}
\def\xsubset{{\,\subset\,}}
\def\xtimes{{\,\times\,}}
\def\xeof{{\,\in\,}}
\def\eq{{\,=\,}}
\def\xto{{\,\to\,}}
\def\eqdef{\,:=\,}
\def\:{{\,:\,}}
\def\ele{{\,\in\,}}
\def\exists{\mbox{ exists }}
\def\and{\mbox{ and }}
\def\for{\mbox{ for }}
\def\fa{\mbox{ for all }}
\def\fe{\mbox{ for each }}
\def\qfa{\quad \mbox{ for all }}
\def\df{\,\dotfill}
\def\iff{if and only if }
\def\wrt{with respect to}

\makeatletter
\newcommand \Df {\leavevmode \cleaders \hb@xt@ .70em{\hss .\hss }\hfill \kern \z@}
\makeatother

\def\sql{\mathfrak{a}}               % sesquilinear form a
\def\DF{B}                           % Dirichlet form a
%------------------------------------------------------------------------------%
% number sets and special elements                                             %
%------------------------------------------------------------------------------%
\def\P{\mathbb{H}}                   % half plane of $\C$
\def\J{I}                            % interval I
\def\N{{\mathbb{N}}}                 % unsigned integers
\def\Q{{\mathbb{Q}}}                 % rational integers
\def\Z{{\mathbb{Z}}}                 % integers
\def\R{{\mathbb{R}}}                 % real numbers
\def\Rn{{\mathbb{R}}^n}              % real numbers in Rn
\def\Rd{{\mathbb{R}}^d}              % real numbers in Rn
\def\C{{\mathbb{C}}}                 % complex numbers
\def\F{{\mathbb{K}}}                 % arbitrary field
\def\Torus{{\mathbb{T}}}             % complex circle line
\def\T{\tau}                         % upper bound of interval (0,t)
\def\sign#1{\operatorname{sign}(#1)} % sign of a number

%------------------------------------------------------------------------------%
% set theory                                                                   %
%------------------------------------------------------------------------------%
\def\G{G}                            % domain $G$
\def\clG{\overline{\G}}              % closure of $G$
\def\bndG{\partial \G}               % boundary of $G$
\def\setA{\mathcal{A}}               % set A
\def\setB{\mathcal{B}}               % set B
\def\setC{\mathcal{C}}               % set C
\def\setM{M}                         % set M
\def\setN{N}                         % set N
\def\setS{\mathcal{S}}               % set S
\def\famF{\mathfrak{F}}              % family of sets
\def\indI{\mathcal{I}}               % index set I
\def\pot#1{\mathfrak{P}(#1)}         % power set

\def\relR{\mathbf{R}}                % relation R
\def\mapf{f}                         % mapping f
\def\mapg{g}                         % mapping g

\def\set#1#2{\{ \, #1 \,: \, #2 \, \}}
\def\tensor#1#2{#1 \otimes #2}
\def\Tens{\oplus}

\def\cal#1{{\mathcal{#1}}}
\def\aut#1{\operatorname{Aut}(#1)}

\def\cross#1#2{#1 {\,\times\,} #2}
\def\multicross#1#2{#1 \times  \ldots \times  #2}
\def\multitens#1#2 {#1 \otimes \ldots \otimes #2}

%------------------------------------------------------------------------------%
% group theory                                                                 %
%------------------------------------------------------------------------------%
\def\1{{\mathsf{1}}}                % unit element of a monoid
\def\0{{\mathsf{0}}}                % unit element of an Abelian monoid
\def\kernel#1{\mathfrak{N}(#1)}     % kernel of a homomorphism
\def\cokernel#1{\mathfrak{C}{(#1)}} % cokernel of a homomorphism
\def\Hom#1{\operatorname{Hom}(#1)}

%------------------------------------------------------------------------------%
% linear algebra                                                               %
%------------------------------------------------------------------------------%
\def\vecV{\mathit{V}}               % vector space V
\def\vecH{\mathit{H}}               % vector space H
\def\vecW{\mathit{W}}               % vector space W
\def\vecU{\mathit{U}}               % subspace U
\def\convC{\mathcal{C}}             % convex set C
\def\basH{\mathfrak{B}}             % Hamel base B

\def\LO#1{\mathscr{L}(#1) }         % algebra of linear transformations

\def\scalar#1#2{ \langle #1,#2 \rangle }
\def\trace#1{\operatorname{tr}(#1)}
\def\dim#1{{\operatorname{dim} \, #1 }}
\def\codim#1{{\operatorname{codim} \, #1 }}
\def\dim{\operatorname{dim}}
\def\codim{\operatorname{codim}}
\def\ind{\operatorname{ind}}

\def\genG{\cal{G}}

\def\algA{\cal{A}}
\def\algB{\cal{B}}
\def\algU{\cal{U}}

\def\idI{\cal{I}}

%------------------------------------------------------------------------------%
% topology                                                                     %
%------------------------------------------------------------------------------%
\def\compK{{K}}          % compact set C
\def\openO{\mathcal{O}}  % open set O
\def\openU{\mathcal{U}}  % open set U
\def\U{\mathcal{U}}      % neighbourhood U
\def\V{\mathcal{V}}      % neighbourhood V
\def\d{\mathit{d}}       % metric function d

\def\interior#1{\mathring{#1}}
\def\closure#1{{\overline{#1}}}

\def\topT{\mathcal{T}}   % topology T
\def\topX{X}             % topological space X
\def\topY{Y}             % topological space Y
\def\topZ{Z}             % topological space Z
\def\metrS{M}            % metric space M

\def\grpG{G}             % topological group G
\def\grpA{A}             % topological group A
\def\grpB{B}             % topological group B
\def\grpC{C}             % topological group C
\def\grpH{H}             % topological group H
\def\grpU{U}             % topological group U
\def\grpV{V}             % topological group V
\def\topG{G}             % topological group G
\def\topH{H}             % topological group H
\def\topU{U}             % topological subgroup U

\def\subS{\cal{S}}
\def\riR{R}              % ring R

%------------------------------------------------------------------------------%
% calculus                                                                     %
%------------------------------------------------------------------------------%
\def\supp#1{\operatorname{supp}(#1)}
\def\singsupp#1{\operatorname{sing \, supp}(#1)}

\def\re{\operatorname{Re}}
\def\im{\operatorname{Im}}
\def\grad{\operatorname{grad}}

%\def\abs#1 {\left|#1\right|}
\def\abs#1 {\lvert #1 \rvert}
%\def\abs#1 {\left|\,#1\,\right|}

%------------------------------------------------------------------------------%
% measure theory                                                               %
%------------------------------------------------------------------------------%
\def\salgA{\Sigma}               % sigma-algebra S
\def\salgB{\cal{B}}              % sigma-algebra B
\def\Borel#1{\cal{B}_{#1}}       % Borel-algebra 
\def\LB{\lambda}                 % Lebesgue-Borel measure

%------------------------------------------------------------------------------%
% function spaces                                                              %
%------------------------------------------------------------------------------%
\def\Lp{L^p(\G)}                 % spaces L_p
\def\Lq{L^q(\G)}                 % spaces L_q
\def\Lh{L^2(\G)}                 % spaces L_2
\def\Li{L^{\infty}(\G)}          % spaces L_infty
\def\Lc{L_{loc}^1(\G)}           % spaces L_1^{loc}
\def\Ck{C^k(\G)}                 % spaces C_k

\def\BV#1{BV(#1) }               % set of functions of bounded variation
\def\AC#1{AC(#1) }               % set of absolutely continuous functions

\def\MP#1{\mathscr{M}^+(#1) }    % set of positive measures
\def\MS#1{\mathscr{M}^-(#1) }    % vector space of finite signed measures
\def\MC#1{\mathscr{M}(#1) }      % vector space of complex measures
\def\MCR#1{\mathscr{M}_{r}(#1) } % vector space of regular complex measures

\def\SobW{W^{k,p}(\G)}                % Sobolev space W
\def\SobO{\mathring{W}^{k,p}(\G)}     % Sobolev space W
%\def\WN#1#2{\mathring{W}^{#1}(#2) }   % Sobolev space W
\def\WN#1#2{W_0^{#1}(#2) }   % Sobolev space W
%\def\HN#1#2{\mathring{H}^{#1}(#2) }   % Sobolev space H
\def\HN#1#2{H_0^{#1}(#2) }   % Sobolev space H
%------------------------------------------------------------------------------%
% functional analysis                                                          %
%------------------------------------------------------------------------------%
\def\snorm#1{\lVert #1 \rVert}
\newcommand{\norm}[1]{\left\lVert #1 \right\rVert}

\def\topV{V}                       % topological vector space V
\def\topW{W}                       % topological vector space W
\def\normS{X}                      % normed space

\def\banX{\mathit{X}}              % Banach space X
\def\banY{\mathit{Y}}              % Banach space Y
\def\banZ{\mathit{Z}}              % Banach space Z
\def\dbanX{\mathit{X^*}}           % dual space of Banach space X
\def\dbanY{\mathit{Y^*}}           % dual space of Banach space Y
\def\dbanZ{\mathit{Z^*}}           % dual space of Banach space Z

\def\banA{\mathit{A}}              % Banach algebra A
\def\banB{\mathit{B}}              % Banach algebra B
\def\banC{\mathit{C}}              % C*-algebra C


\def\domain#1{\mathfrak{D}(#1)}    % domain
\def\range#1{\mathfrak{R}(#1)}     % range

\def\isoJ{\mathfrak{J}}            % isometry J
\def\repA{{\cal{A}_{\sql}}}        % representation operator A
\def\linS{{S}}                     % linear operator S
\def\linG{{G}}                     % linear operator G
\def\linL{{L}}                     % linear operator L
\def\linT{{T}}                     % linear operator A
\def\linA{\mathit{A}}              % linear operator A
\def\linB{{B}}                     % linear operator B
\def\linM{{M}}                     % linear operator M
\def\linU{{U}}                     % linear operator U
\def\linI{{I}}                     % linear operator I
\def\A{\cal{A}}                    % linear differential operator A
\def\BDO{\cal{B}}                  % linear boundary differential operator B
\def\linU{{U}}                     % unitary operator U
\def\linP{{P}}                     % projection P
\def\compA{k}                      % compact operator k
\def\linFR{f}                      % finite rank operator f
\def\fredA{A}                      % Fredholm operator F

\def\HAM{\cal{H}}                  % Hamilton operator
\def\PA{\partial^{\alpha}}         % elementary differential operator
\def\DO{\cal{A}}                   % linear differential operator
\def\BC{\cal{B}}                   % boundary operator
\def\SLO{\cal{L}_{p,q}}            % Sturm-Liouville operator
\def\MO{\cal{L}_{M}}               % momentum operator

\def\HS#1{S(#1)           }        % algebra of Hilbert-Schmidt operators
\def\FR#1{F(#1)           }        % algebra of finite rank operators
\def\TC#1{N(#1)           }        % algebra of trace class operators
\def\BU#1{\mathscr{U}(#1) }        % algebra of unitary linear operators
\def\BO#1{\mathscr{L}(#1) }        % algebra of bounded linear operators
\def\CO#1{\mathscr{C}(#1) }        % algebra of compact linear operators
\def\CL#1{\mathscr{C}(#1) }        % algebra of compact linear operators
\def\FO#1{\Phi(#1) }               % set of Fredholm operators
\def\FK#1#2{\Phi_{#2}(#1) }        % set of Fredholm operators with index k

\def\hilH{\mathit{H}}              % Hilbert space H
\def\clU{\mathit{U}}               % closed subspace U

\def\orthO{\mathfrak{O}}           % orthonormal system O
\def\orthS{\mathfrak{S}}           % orthonormal system S
\def\basB{\mathfrak{B}}            % basis B of a Hilbert space

\def\GL#1 {\mathbf{GL}(#1)}        % linear group
\def\OG#1 {\mathbf{O}(#1)}         % linear group
\def\SOG#1 {\mathbf{SO}(#1)}       % linear group

%------------------------------------------------------------------------------%
% distribution theory                                                          %
%------------------------------------------------------------------------------%
\def\disF{F}                  % distribution F
\def\disG{G}                  % distribution G
\def\disH{H}                  % distribution H
\def\disU{U}                  % distribution U
\def\disT{T}                  % distribution T
\def\SF{C^{\infty}(\G)}       % space of smooth functions
\def\TF{C_c^{\infty}(\G)}     % space of compactly supported smooth functions
\def\TFRd{C_c^{\infty}(\R^d)} % space of test functions
\def\SRd{\mathscr{S}(\R^d)}   % Schwartz space
\def\SR1{\mathscr{S}(\R)}     % Schwartz space
\def\B1#1{C^1(#1) }           % space of continuously differentiable functions
\def\Bm#1{C^m(#1) }           % space of continuously differentiable functions
\def\Ck{C^k(\G)}
\def\TD{\mathscr{S}'(\R^d)}   % space of temperated distributions
\def\E{\cal{E}}               % space of distributions with compact support
\def\D{\mathscr{D}}           % D
\def\FT{\mathcal{F}}          % Fourier transformation F
\def\FT{\mathscr{F}}          % Fourier transformation F

%------------------------------------------------------------------------------%
% other definitions                                                            %
%------------------------------------------------------------------------------%
\def\Proof{\textsc{Proof}}
\def\FRECHET{Fr\'{e}chet}
\def\ARZELA{Arzel\'{a}}
\def\GARDING{G\r{a}rding}
\def\HOELDER{H\"older}
\def\SCHROEDINGER{Schr\"odinger}
\def\JOERGENS{J\"orgens}
\def\WUEST{W\"uest}
\def\KAEHLER{K\"aehler}
\def\HOEFLINGER{H\"oflinger}
\def\POINCARE{Poincar\'{e}}
\def\FA{Functional Analysis}

\def\Author   {K.\,{\HOEFLINGER}}
\def\AuthorUC {K.\,HOEFLINGER}
\def\Version  {1.0}
\def\Email    {khoeflinger@t-online.de}


%\renewcommand\thesection{\Roman{section}}

\def\Title{Spectral Theory of Partial Differential Operators}

\titleformat{\section}
{\normalfont \bfseries \filcenter}
{\thesection.\;}{0mm}{}

\titleformat{\subsection}%[runin]
{\normalfont \itshape}
{}{}{}

\titlespacing*{\section}   {0pt}{12pt}{4pt}
\titlespacing*{\subsection}{0pt}{6pt}{0pt}

%------------------------------------------------------------------------------%
%                                BEGIN OF DOCUMENT                             %
%------------------------------------------------------------------------------%
\begin{document}
\thispagestyle{empty}

\centerline{\large{\textbf{\Title}}}

\vspace{4pt}
\centerline{\footnotesize{{\textsc{Joachim Weidmann}}}}
\vspace{6pt}

\fancyhead[C] {\small{\textsc{\Title}}}
\fancyhead[OL]{\small{\thepage}}
\fancyhead[ER]{\small{\thepage}}

\myfootnote
{
  Reprint of a lecture published
  in \textit{Spectral Theory and Differential Equations},
  Lecture Notes in Mathematics 448, Springer Verlag,
  pp. 71-111.
}

%------------------------------------------------------------------------------%
%\section*{Introduction}                                                       %
%------------------------------------------------------------------------------%
In this paper we review some methods of \textit{perturbation theory}
which are especially useful for the study of
\textit{spectral properties of self-adjoint differential operators}.

\subparagraph*{}
First we describe self-adjoint differential operators with constant and
not constant coefficients in $L_2(\R^m)$ or $L_2(G)$,
where $G$ is an open subset of $\R^m$ (sections 1 and 2).
In section 3 we give some abstract perturbation results concerning
perturbations which are relatively compact with respect to the unperturbed
operator or with respect to the square of the unperturbed operator.
Roughly speaking the results say that the essential spectrum
(and frequently the singular sequences) are preserved.
In section 4 we show how these results are applicable
to differential operators,
especially to operator of the \textit{{\SCHROEDINGER}}
or \textit{Dirac type}.
We study perturbations which are \textit{small at infinity}.
It turns out that these perturbations are very well adapted to this situation.
In section 5 finally we show that similar conditions are also useful
in \textit{scattering theory}
(and therefore in the perturbation theory
of the absolutely continuous spectrum).

\paragraph*{}
We introduce some notations.
With $L_2(G)$ we denote the Hilbert space of square integrable complex-valued
functions.
By $L_{2,M}(G) := \{L_2(G)\}^M \eq L_2(G,\C^M)$ we denote
the corresponding space of $\C^M$-valued functions.
The inner products and norms are in both cases denoted
by $\scalar{\,\cdot\,}{\,\cdot\,}$ and $\norm{\,\cdot\,}$,
respectively.
The additional index $M$ is used in the same sense in other cases:
$C_{0,M}^{\infty}$, $W_{2,M}^r$ etc.

\subparagraph*{}
The \textit{Fourier transform} in $L_2(\R^m)$
(and in the same way in $L_{2,M}(\R^m)$) is defined by
\[
F\,f(x) \= (2\pi)^{-m/2} \lim_{N \to \infty}
\int_{\abs{y} \leq N} \,e^{-ixy} f(y) \,dy;
\]
here $xy$ stands for the inner product of $x,y \in \R^m$.
With $b_j = i \frac{\partial}{\partial x_j}$ we then have
\[
F \partial_j f = -x_j F\,f,
\quad
F\,x_jf = \partial_j F\,f,
\]
where $x_j$ is the $j$-the component of $x$.

%------------------------------------------------------------------------------%
\section{Differential operators with constant coefficients}                    %
%------------------------------------------------------------------------------%
The simplest differential operators are of course those
with \textit{constant coefficients}.
Such an operator $T$ can be written in the form
\[
T = P(\partial) = P(b_1,\ldots,b_m),
\]
where $P$ is a polynomial in $m$ variables,
\[
b = (\partial_1,\ldots,\partial_m), \quad b_j = i \frac{\partial}{\partial x_j}
\]
(the imaginary factor $i$ is included to make $b_j$
\textit{formally self-adjoint};
this means, for $u,v \in C_0^{\infty}(\R^m)$ we have
$\scalar{b_j u}{v} = \scalar{u}{b_j v}$.
A quite natural domain of definition for such a differential operator
is $C_0^{\infty}(\R^m)$,
which is dense in $L_2(\R^m)$.

\subparagraph*{}
For $u,v \in C_0^{\infty}(\R^m)$ we have
\[
\scalar{P(\partial) u}{v} = \scalar{u}{P^*(\partial) v},
\]
where $P^*$ is the polynomial with coefficients
which are conjugate to the coefficients of $P$.
Therefore $T_0$,
the restriction of $T$ to $\domain{T_0} = C_0^{\infty}(\R^m)$,
is symmetric if the coefficients of $P$ are real.

\subparagraph*{}
By means of the Fourier transform $F$ the operator $T_0$ is transformed to
\[
M_{P,0} = F^{-1} \, T_0 \, F,
\]
where $M_{P,0}$ is the restriction of the multiplication by $P$ to the domain
$\domain{M_{P,0}} = F^{-1} C_0^{\infty}(\R^m)$,
which is also dense in $L_2(\R^m)$.
From that it is clear that $M_{P,0}$,
and therefore $T_0$, are closable.
$D{\overline{M_{P,0}}}$,
and therefore $\domain{\overline{T_0}}$,
obviously contain the Schwartz space $S(\R^m)$;
this immediately implies that $\overline{M_{P,0}} = M_P$,
the maximal operator of multiplication by $P$, {\ie}
\[
\domain{M_P} = \{f \in L_2(\R^m): P\,f \in L_2(\R^m) \},
\]
\[
M_P\,f = P\,f \for f \in \domain{M_p}.
\]
Since $M_P$ is self-adjoint if and only if $P$ is real valued
({\ie} the coefficients of $P$ are real),
the operator $T = \overline{T_0}$ is \textit{self-adjoint}
if and only if $P$ has real coefficients.
This implies that $T_0$ is \textit{essentially self-adjoint}
in this case, and
\[
T = T_0^* = F \, M_p \, F^{-1}.
\]

\subparagraph*{}
In general $T$ is a \textit{normal} operator, {\ie}
\[
\domain{T} = \domain{T^*},
\quad
\norm{T^* f} = \norm{T f} \for f \in \domain{T}.
\]
$T^*$ is the operator generated by $P^*$ in the same sense as $T$ was generated
by $P$.
(We could say: $T_0$ is essentially normal,
since $\overline{T_0}$ is normal).

\subparagraph*{}
The domain of $T$ apparently is given by
\[
\domain{T} = F\,\domain{M_P} = \{ f \in L_2(\R^m): P\,F^{-1}f \in L_2(\R^m)\}.
\]
From $T = T_0^*$ and the definition of the derivatives for distributions
it follows that $f \in \domain{T}$ if and only if $P(\partial)f \in L_2(\R^m)$,
$T\,f = P(\partial) f$.
(Here the derivatives are taken in the sense of distributions.
This means: take the regular distribution $D_f$ which is generated by $f$,
calculate $P(\partial) D_f$ in the sense of distributions;
this is again a regular distribution generated by an $L_2$-function $g$;
define $P(\partial)f = g$.)
This is another useful description of $\domain{T}$ and $T$.

\subparagraph*{}
Let $P_1$ and $P_2$ be polynomials with
\[
\abs{P_1(\xi)} \leq C_1 + C_2 \abs{P_2(\xi)} ,
\quad \xi \in \R^m
\]
for suitable constants $C_1$ and $C_2$.
Let further $T_1$ and $T_2$ be the operators generated by $P_1$ and $P_2$,
respectively.
It follows that $\domain{T_1} \subset \domain{T_2}$
and for a suitable constant $C$
\[
\norm{T_1 f} \leq C \norm{f} + \norm{T_2 f}.
\]
For example, if
\[
T := -\Delta = \sum_{j=1}^n \partial_j^2,
\]
then every derivative (in the sense of distributions) of $f \in \domain{T}$
up to order $2$ lies in $L_2(\R^m)$.
If,
on the other hand,
$T = \partial_1 \partial_2$,
then $\partial_1 \partial_2 f$ is (in general) the only derivative
of $f$ which lies in $L_2(\R^m)$ for $f \in \domain{T}$.

\subparagraph*{}
If $T$ is an \textit{elliptic} operator of order $r$,
{\ie} $P$ is a polynomial of order $r$ and
\[
C_1 + \abs{P(\xi)} \geq C_2 \abs{\xi} ^r,
\quad \xi \in \R^m
\]
for some constants $C_1$ and $C_2$,
then $\domain{T}$ is equal to the well-known \textit{Sobolev space}
\[
W_2^r(\R^m) := \{ f \in L_2(\R^m): \abs{x} ^r F^{-1}f \in L_2(\R^m) \}
\]
\[
\= \{ f \in L_2(\R^m): \abs{x} ^r Ff \in L_2(\R^m) \}.
\]
This definition is meaningful for all $r \in \R$,
and we shall use it later for all positive $r$.
For $r=0$ we have $W_2^r(\R^m) = L_2(\R^m)$.
From the above considerations it follows that
$C_0^{\infty}(\R^m)$ is \textit{dense} in $W_2^r(\R^m)$
with respect to the \textit{Sobolev norm}
\[
\norm{f}_r := \norm{f}_{W_2^r} = \norm{(1+\abs{\cdot} )^r) Ff}.
\]
With this norm $W_2^r(\R^m)$ becomes a Hilbert space.
One of the most important facts about these Sobolev spaces is given
by the following theorem.

\begin{theorem}
Let $0 \leq s < r$.
Then $W_2^r(\R^m) \subset W_2^s(\R^m)$ and the embedding is continuous.
For every $\epsilon > 0$ there is a $C \geq 0$ such that
\[
\norm{f}_s \leq \epsilon \norm{f}_r + C \norm{f}
\mbox{ for every } f \in W_2^r(\R^m).
\]
For every $\phi \in C_0^{\infty}(\R^m)$ and $f \in W_2^r(\R^m)$ we have
$\phi f \in W_2^r(\R^m)$ and the map
\[
J_\phi: W_2^r(\R^m) \to W_2^s(\R^m), \quad f \mapsto \phi f
\]
is compact.
\end{theorem}

\subparagraph*{}
The first and second statement are immediate consequences of the
definition of the norms in $W_2^r(\R^m)$.
The remaining results follow easily from the fact that the multiplication
by $\phi$ is represented by a convolution with the rapidly decreasing
function $(2\pi)^{-m/2} F^{-1} \phi$ in the Fourier representation.

\subparagraph*{}
Sometimes local Sobolev space are used.
We define
\[
W^r_{2,loc}(\R^m) := \{ f \in L_{2,loc}(\R^m):
\phi f \in W_2^r(\R^m)
\mbox{ for } \phi \in C_0^{\infty}(\R^m) \}.
\]
It is possible to define an inductive limit topology such that
$W^r_{2,loc}(\R^m)$ becomes a (locally convex) {\FRECHET} space;
but we do not need such a topological structure.

\subparagraph*{}
If $P$ does not satisfy the ellipticity condition,
then one could define a "generalized Sobolev space" by
\[
W_2(\R^m,P) := \{ f \in L_2(\R^m): PF^{-1}f \in L_2(\R^m) \}.
\]
It is easy to prove some quite similar results for such spaces,
but we do not go into details.

\subparagraph*{}
All questions concerning \textit{spectral properties} of the operators
considered above are now almost trivial,
since they are \textit{unitary equivalent} to multiplication operators
(in fact to multiplication by polynomials).
By means of the unitary equivalence we see that the spectrum of $T$
is equal to the range of $P$
(remember that the range of a polynomial is always closed).
Therefore,
the spectrum of a self-adjoint differential operator $T$ is either
a half plane (if $T$ is semi-bounded),
or the entire real line (if $T$ is not semi-bounded).
Since $\lambda$ is an eigenvalue of the operator of multiplication by
a function $t$ if and only if $t^{-1}(\lambda)$ has positive measure,
if follows that our operator $T$ have \textit{no eigenvalues}.

\paragraph*{}
The \textit{spectral resolution} $\hat{E}$ of the self-adjoint operator
$M_P$ is given by
\[
(\hat{E}(b) - \hat{E}a) f \= 
\chi_{\{x \in \R^m: a \leq P(x) \leq b\}} f.
\]
From this we can see that the spectra of $M_P$ and $T$ are
\textit{absolutely continuous},
which means that the measure on $\R$,
which is generated by
\[
\mu([a,b)) \= \scalar{f}{(\hat{E}(b) - \hat{E}a) f}
\]
is absolutely continuous with respect to the Lebesgue measure
for every $f \in L_2(\R^m)$.
The spectral projection $\hat{E}(b)-\hat{E}(a)$ of $T$
is given by the convolution with
\[
(2\pi)^{-m/2} F \chi_{\{x \in \R^m: a \leq P(x) \leq b\}}
\]
(which lies in $L_2(\R^m)$.

\paragraph*{}
The \textit{resolvent} of $M_P$ is the multiplication operator
\[
(\lambda - M_P)^{-1} f(x) \= (\lambda - P(x))^{-1} f(x)
\]
for $\lambda$ not in the range of $P$.
Therefore the resolvent of $T$ can be given as a convolution operator,
but in general this is possible only in the sense of distributions.

\subparagraph*{}
For the very important operator $T \eq -\Delta$ in $\R^3$
we have $\sigma(T) \eq [0,\infty)$ and 
\[
(\lambda - T)^{-1} f(x) \= \frac{-1}{4\pi}
\int \abs{x-y} ^{-1} \, \exp [i \sqrt{\lambda} \abs{x-y} ] \, f(y) \, dy,
\]
where $\im \sqrt{\lambda} > 0$ for $\lambda \neq [0,\infty)$.
It is also possible to calculate the \textit{unitary group} $exp(-itT)$
by means of the Fourier transform,
since the unitary group of $M_P$ is multiplication by $exp(-itP)$.
We get for $T \eq -\Delta$ and $m \in \N$
\[
(e^{-itT}f)(x) = \lim_{N \to \infty} (4\pi it)^{-m/2}
\int_{\abs{x} \leq N} e^{-\abs{x-y} /4it} \, f(y)\,dy.
\]

\subparagraph*{}
These things are only a little more complicated for differential operators
in the Hilbert spaces of $C^M$-valued functions
$L_{2,M}(\R^M) := L_2(\R^m)^M = L_2(\R^m,\C^M)$.
In this case we assume that $P$ is a matrix valued function on $\R^m$,
where all the entries are polynomials.
Again by means of the Fourier transform on $L_{2,M}(\R^m)$,
we get
\[
T_0 \= F \, M_{P,0} \, F^{-1},
\]
where now $M_{P,0}$ is the operator of (matrix-vector) multiplication by $P$.
By the same technique as above we prove that $T_0$ is closable and
that $T_0^*$ is unitarily equivalent to the maximal multiplication operator
$M_{P^*}$,
where $P^*(x)$ is the matrix which is hermitean conjugate to $P(x)$.
It follows that $T := \overline{T_0}$ is essentially self-adjoint
if and only if $P(x) = P^*(x)$; then $T = T_0^*$.

\subparagraph*{}
one of the most important examples of this kind is the \textit{Dirac operator}
for a free electron.
In this case we have $m=3$, $M=4$ and
\[
P(x) \= c(y_1 \alpha_1 + y_2 \alpha_2 + y_3 \alpha_3) + \mu c^2 \beta,
\]
$\mu$ and $c$ are positive constants ($\mu$ the mass of the electron,
$c$ the velocity of the light),
\[
a_j \=
\begin{pmatrix}
0        & \sigma_j \\
\sigma_j & 0
\end{pmatrix}
\quad (j=1,2,3),
\beta \=
\begin{pmatrix}
\;1 & \;0 & \;0 & \;0 \\
\;0 & \;1 & \;0 & \;0 \\
\;0 & \;0 &  -1 & \;0 \\
\;0 & \;0 & \;0 &  -1
\end{pmatrix}
\]
and
\[
      \sigma_1 \= \begin{pmatrix} 0 & \; 1 \\ 1 & \; 0 \end{pmatrix},
\quad \sigma_2 \= \begin{pmatrix} 0 &   -i \\ 1 & \; 0 \end{pmatrix},
\quad \sigma_3 \= \begin{pmatrix} 1 & \; 0 \\ 0 &   -1 \end{pmatrix}.
\]
Since $P^*(x) = P(x)$ the operator $T_0$ is essentially self-adjoint;
$T := \overline{T_0}$ is self-adjoint with the domain
$W^1_{2,M}(\R^m) = W^1_2(\R^m)^M$.

\subparagraph*{}
The eigenvalues of the matrix $P(x)$ are given by $h_j(x)$ with
\[
h_1(x) = h_2(x) = -h_3(x) = -h_4(x)
= \mu c^2 (1 + \frac{\abs{x} ^2}{\mu^2 c^2})^{1/2}.
\]
These eigenvalues and a suitable choice of corresponding eigenvectors
depend smoothly on $x$ (see [25], section 4).
Therefore $T$ is unitarily equivalent to the multiplication operator $M_H$,
where
\[
H(x) \=
\begin{pmatrix}
h_1(x) & 0      & 0      & 0      \\
0      & h_2(x) & 0      & 0      \\
0      & 0      & h_3(x) & 0      \\
0      & 0      & 0      & h_4(x)
\end{pmatrix}.
\]
from which we can easily see that the \textit{spectrum} is
\[
\sigma(T) = (-\infty, -\mu c^2] \, \cup \, [\mu c^2,\infty).
\]
It is also not difficult to show
that this spectrum is \textit{absolutely continuous}.

\paragraph*{}
Finally we should mention that the differential operators are just
a subclass of those operators
which are \textit{Fourier transforms of multiplication operators}
(in the same sense as differential operators are Fourier transforms
of multiplications by polynomials).
This more general class of operators has many properties
which are the same as for differential operators.
If, for example,
the function $P \: \R^m \xto \R$ is smooth and $\grad P \neq 0$
almost everywhere,
then $M_P$ and therefore $T$) has an absolutely continuous spectrum
and there exists at least some part of a reasonable scattering theory,
e.\,g. [24,25].
One great difference lies in the fact
that differential operators are \textit{local operators}
({\ie} $\supp{T \phi} \subset \supp{\phi}$,
but these more general operators are not local.

%------------------------------------------------------------------------------e
\section{Perturbed operators}                                                  %
%------------------------------------------------------------------------------%
Most operators which appear in applications do not have constant coefficients,
but in many cases the terms of hoghest order have constant
(or asymptotically constant) coefficients.
Therefore it seems reasonable to consider these operators as
\textit{perturbed operators},
the \textit{unperturbed operator} being one with constant coefficients.
Therefore all the basic problems of spectral theory can be studied
from the view point of \textit{perturbation theory}:
assume that we know the unperturbed operator $T_0$ very well;
what can we say about the perturbed operator.
Almost everything we are goint to describe in this paper is
perturbation theory in some sense.

\paragraph*{}
The first problem is,
to give a reasonable definition of the perturbed operator,
{\ie} to find a domain of definition such that
this operator is self-adjoint or, at least, essentially self-adjoint.
One of the most useful (abstract) results was given by Rellich and Kato
(see [12], V.4).

\subparagraph*{}
Let $A$ and $B$ be operators in a Hilbert space $H$.
We say that $B$ is \textit{$A$-bounded},
if $\domain{A} \subset \domain{B}$and
if there exist non-negative numbers $a$ and $b$ such that
\[
\norm{Bx} \leq a \norm{Ax} + b\norm{x}, \quad x \in \domain{A}.
\]
The infimum of all possible numbers $a$ is called the \textit{$A$-bound} of $B$.

\begin{theorem}
Let $A$ be a (essentially) self-ajoint operator in a Hilbert space $H$
and assume that $B$ is symmetric and $A$-bounded with $A$-bound less than $1$.
The $A+B$ is (essentially) self-adjoint and
$\overline{A+B} \overline{A} + \overline{B}$
(this implies $\domain{\overline{A+B}} = \domain{A}$.
\end{theorem}

\subparagraph*{}
It turns out,
that very often on eproves the $A$-boundedness with relative bound $<1$
of a perturbation $B$ by proving that it is $A$-bounde woth $A-bound$ $0$.
This is so,
if one considers differential operators where the coefficients are taken
from a suitable \textit{Stummel class}.

\subparagraph*{}
For a real number $\rho$ we say that a function $q \:\R^m \xto \C$ belongs to
the \textit{local Stummel class} $M_{\rho,loc}(\R^m)$, if
\[
M_{q,\rho}(x) \eqdef
\begin{cases}[1.4]
\quad (\int_{\abs{x-y} \leq 1} \; \abs{q(y)} ^2 \abs{x-y} ^{\rho-m} \, dy)^{1/2}
& \for p < m,\\
\quad (\int_{\abs{x-y} \leq 1} \; \abs{q(y)} ^2 \, dy)^{1/2}
& \for p \geq m\\
\end{cases}
\]
exists for every $x \in \R^m$ and if $M_{q,\rho}(\cdot)$ is locally bounded.
For $\rho \geq m$ we have $M_{\rho,loc}(\R^m) \eq L_{2,loc}(\R^m)$.
We say that
$q$ belongs to the (global) \textit{Stummel class} $M_{\rho}(\R^m)$,
if
\[
M_{q,p} \eqdef \sup_{x \in \R^m} M_{q,p}(x) \, < \, \infty. 
\]
We have obviously $M_{\rho,loc} \subset M_{\tau,loc}$, if $\rho \leq \tau$.
With this notion the following theorem holds.

% xxx Theorem 2.2.

\paragraph*{}
A simple consequence of this theorem is

% xxx Theorem 2.3.

% xxx 2.4 example
% xxx 2.5 example

\paragraph*{}
By means of very specific considerations it can be shown
(see Kato [12], V.5) that,
in order to have self-adjointness of the perturbed Dirac operator
on the same domain as $T_0$,
$q$ may also contain sufficiently small Coulomb terms;
but in this case the perturbation is not relatively compact with respect to $T_0$
(relative bound $> 0$).

\paragraph*{}
Much more general conditions are possible
if we do only wish that $T = T_0 + V$ is essentially self-adjoint
on $C_0^{\infty}(\R^m)$
(but not,
that the self-adjoint extension of $T$ has the same domain as that of $T_0$).
For the {\SCHROEDINGER} operator (Example 2.4) the first general result
of this kind was given by Ikebe and Kato [8]
(notice that they considered operators with variable highest order coefficients).
Recently this result has been generalized by Simon [21,22,23] and Kato [13].
The latest, and probably most general result was recently given by Simader [20].
The following theorem combines the results of Ikebe-Kato and Simader.

% xxx Theorem 2.6

\paragraph*{}
One can prove similar results for the Dirac operator.
One of the most general results in this direction was given by {\JOERGENS} [10].

% xxx Theorem 2.7

\paragraph*{}
For many application it is also important to study differential operators
in $L_2(G)$, where $G$ is an open subset of $\R^m$.
Conditions for the essential self-adjointness on $C_0^{\infty}(\R^m)$
of operators of the form considered in 2.4 (the operators of {\SCHROEDINGER} type
and even for a more general class of operators) have been given
by {\JOERGENS} [9].
In this case,
similar to the Ikebe-Kato result,
the domain of the self-adjoint extensions consist of all functions
$u \in L_2(G) \cap W_{2,loc}^2(G)$ with $T\,u \in L_2(G)$;
here $T\,u$ is calculated in the sense of distributions;
the local \textit{Sobolev space} $W_{2,loc}^r(G)$
is the set of all $u \in L_{2,loc}(G)$
such that $\phi u \in W_2^r(\R^m)$ for every $\phi \in C_0^{\infty}(G)$.
No result of this kind seems to be known for a more general class of differential
operators; actually it seems very difficult to guess how such a result might look!

%------------------------------------------------------------------------------%
\section{The invariance of the essential spectrum; abstract results}           %
%------------------------------------------------------------------------------%
The \textit{essential spectrum} $\sigma_e(A)$ of a self-adjoint operator $A$
in a Hilbert space $H$ is defined to be the set consisting of all limit points
of the spectrum and the eigenvalues of infinite multiplicity.
It is known that a real number $\lambda$ belongs to $\sigma_e(A)$ \iff one of
the following conditions is satisfied:

\begin{itemize}[leftmargin=6mm]
\item[3.1]
For every $\epsilon > 0$
the projection $E(\lambda+\epsilon)-E(\lambda-\epsilon)$
is infinite-dimensional, or
\item[3.2]
there exists a sequence $(u_n)$ in $\domain{A}$ such that $u_n \to \0$,
$u_n \neq \0$ and $(A-\lambda)u_n \to \0$;
such a sequence is called a \textit{singular sequence} for $A$ and $\lambda$.
\end{itemize}
An important problem is the \textit{invariance of the essential spectrum}
under perturbations.
In this connection \textit{relative compactness} plays an important role.

\paragraph*{}
Let $A$ and $B$ be operators in $H$.
$B$ is said to be $A$\textit{-compact} if $\domain{A} \subset \domain{B}$
and the map $B \: \domain{A} \xto H$ is compact
(here $\domain{A}$ is equiped with the \textit{graph norm}
$\norm{x}_A := (\norm{x}^2 + \norm{Ax}^2)^{1/2}$.
An equivalent condition is:
from $x_n\to \0$ and $Ax_n\to \0$ it follows that $Bx_n \to \0$
(for details see [11]).

\subparagraph*{}
For a self-adjoint operator $A$ with the spectral resolution $E$ it follows
easily from the von Neumann spectral theorem:
$B$ is $A$-compact if and only if $B$ is $A$-bounded with $A$-bound $0$ and
the restriction of $B$ to $\range{E(N)-E(-N)}$ is compact for every $N \geq 0$.

\paragraph*{}
The following result is well-known.
It follows immediately from the above characterizations
of the essential spectrum and of the relative compactness.

\begin{theorem} % theorem 3.3
Let $A$ be self-adjoint and let $B$ be symmetric and $A$-compact.
Then $\sigma_e(A+b) = \sigma_e(A)$.
\end{theorem}

\subparagraph*{}
For many interesting applications to differential operators the
requirement of $A$-compactness is too restrictive.
Recently Schechter[18] proved that in certain cases $A^2$-compactness is
also sufficient.
His result holds also for non self-adjoint operators in Banach spaces;
the proof makes strong use of Fredholm theory.
For this reason we will give here a simpler proof for self-adjoint operators.

\begin{theorem} % theorem 3.4
Let $A$ be a self-adjoint operator in $H$, $B$ symmetric and $A$-bounded
such that $A+B$ is self-adjoint.
If $B$ is $A^2$-compact,
then $\sigma_e(A+b) = \sigma_e(A)$.
\end{theorem}

% xxx 3.5. Remarks
% xxx Proof of theorem 3.4

\paragraph*{}
So far we have quite general conditions which guarantee the invariance
of the essential spectrum;
in the following section we shall give simple conditions which are sufficient
for $A^2$-compactness if $A$ is a self-adjoint differential operator.

\paragraph*{}
We should note that the conditions of Theorem 3.4 are, in general,
not sufficient to guarentee that $A+B$ is semi-bounded if $A$ is semi-bounded.
As an example consider $A$ in $L_2(\R^m)$ defined by
\[
xxx.
\]
For $B$ we chose the operator defined by
\[
xxx.
\]
Then it is easy to show that $B$ is $A$-bounded and $A^2$-compact,
and $A+B$ is self-adjoint.
But $-\infty$ is a cluster point of eigenvalues of $A+B$.

\subparagraph*{}
A simpler example is as follows:
Let $A$ be a positive self-adjoint operator with a discrete spectrum
and $B = -2A$. -- Such a situation can of course not occur
if $B$ is $A$-bounded with $A$-bound less than $1$.

\paragraph*{}
It can be expected that sharper results hold, if the perturbation os semi-definite.
For example, if $A$ has the spectrum $[\mu,\infty)$ and $B$ nonnegative,
then the spectrum of $A+B$ is contained in $[\mu,\infty)$.
But what happens, if there are holes in the spectrum (or in the essential spectrum)
of $A$, as it occurs, for example, in the case of the Dirac-operator?
A first result of this type is given by Glazman [5],theorem 7$^{bis}$;
this theorem is the same as our Corollary 3.9\,b.

% xxx theorem 3.6

\paragraph*{}
There are some simple corollaries of theorem 3.6.
The proofs are easy and will be left to the reader.

% xxx corollary 3.7
% xxx corollary 3.8
% xxx corollary 3.9
% xxx theorem 3.10

\paragraph*{}
It seems remarkable that this result can be extended not only to $A$-compact
perturbations, but also th some $A^2$-compact perturbations.

% xxx theorem 3.11

%------------------------------------------------------------------------------%
\section{Perturbations which are small at infinity}                            %
%------------------------------------------------------------------------------%
In [11] J\"{o}rgens and the author introduced a general class of perturbations
which are "relatively small at infinity".

% xxx rest of the section

%------------------------------------------------------------------------------%
\section{The absolutely continuous spectrum and scattering theory}             %
%------------------------------------------------------------------------------%
From examples (see for example Aronzajn [1]) it is well known that
we cannot expect a reasonable perturbation theory for the continuous spectrum,
neither for differential operators nor in an abstract setting.

\subparagraph*{}
From scattering theory we know that the situation is much better
with the \textit{absolutely continuous spectrum}.
But, let us first give some definitions.
If $T$ is a self-adjoint operator in a Hilbert space $H$,
then the \textit{absolutely continuous subspace} $H_{ac}$ of $T$
is defined by
\[
H_{ac} = \{ u \in H: \scalar{E(\cdot)u}{u} \mbox{ is absolutely continuous}\}.
\]
The restriction $T_{ac}$ of $T$ to $H_{ac}$ is called the \textit{absolutely
continuous part} of $T$;
the spectrum of $T_{ac}$ is the \textit{absolutely continuous spectrum} of $T$,
$\sigma_{ac}(T)$.
In this section we consider operators $T_0$ and $T_1$;
to all symbols $H_{ac}$, $\sigma_{ac}$, $T_{ac}$ we add the index $0$ or $1$
if they refer to $T_0$ or $T_1$ respectively.

\subparagraph*{}
Let us now assume that $T_0$ has a purely absolutely continuous spectrum
({\ie} $H_{ac,0} = H$);
this is always the case if $T_0$ is a differential operator with constant
coefficients, $T_0 \neq 0$.
Then the \textit{wave operators} $W_{\pm}(T_1,T_0)$ are defined by
\[
W_{\pm}(T_1,T_0) \=
\mbox{s-}\lim_{t \to \pm \infty} \exp(itT_1) \, \exp(-itT_0)
\]
if they exist ("s-lim" means "strong limit").

\subparagraph*{}
If the wave operators $W_{\pm} = W_{\pm}(T_1,T_0)$ exist,
then these are isometries and the ranges $\range{W_{\pm}}$ are subspaces of
$H_{ac,1}$ and $T_0$ is \textit{unitarily equivalent}
to $T_1\mid_{\range{W_{\pm}}}$.
In general this is not very much information,
since we do not know anything about the restriction of $T_1$
to $\range{W_{\pm}}^{\bot}$.
But there are results which guarantee that the ranges of $W_{\pm}$ are equal
to all of $H_{ac,1}$.
One of these results is given by

%-- theorem 5.1.
\begin{theorem}([12], note the end of X.4.5)
If $(\lambda-T_1)^{-n} - (\lambda-T_0)^{-n}$ is in the trace class
for every non-real $\lambda$,
then the wave operators $W_{\pm}(T_1,T_0)$ exist and
$\range{W_+} \eq \range{W_-} \eq H_{ac,1}$.
(The simplest characterization of an operator of the \textit{trace class} is,
that it is a product of two Hilbert-Schmidt operators.)
\end{theorem}

\paragraph*{}
This result has been proved in [26] to the $1$-dimensional case.
There we have proven essentially the following theorem.

%-- theorem 5.2.
\begin{theorem}
Let $T_0$ be a self-adjoint (elliptic) differential operator with constant
coefficients in $L_2(\R^m)$ of order $>2$ and assume that $V$ is a symmetric
operator in $L_2(\R^m)$ such that:
$\domain{T_0} \subset \domain{V}$,
$T_1 = T_0 + V$ is self-adjoint and there exist constants $C \geq 0$ and
$\Theta < -1$ such that for every $r \geq 0$
\[
\norm{V u} \leq C (1+r)^{\Theta} (\norm{u} + \norm{T_0}),
\]
if $u \in \domain{T_0}$ ans $\abs{\supp{u}} \geq r$.
Then the wave operators $W_{\pm}(T_1,T_0)$ exist and we have 
$\range{W_+} \eq \range{W_-} \eq H_{ac,1}$.
\end{theorem}

\subparagraph*{}
So far we are not able to give such a general result for higher dimensions
of for $m \eq 1$ and an operator $T_0$ or order $1$.
But it is not difficult to prove the following extension which covers
"smooth and reapidly decreasing" perturbations.

%-- theorem 5.3.
\begin{theorem}
Let $T_0$ be a self-adjoint elliptic differential operator
with constant coefficients in $L_2(\R^m)$ of order $r \geq 1$.
Let $V$ be a symmetric operator in $L_2(\R^m)$ such that
$\domain{T_0} \subset \domain{V}$, $T_1 = T_0+V$ is self-adjoint and there
exist constants $C \geq 0$ and $\theta < -m$ such that for every $r \geq 0$
\[
\norm{Vu} \leq C (1+r)^{\theta} (\norm{u}+\norm{T_0 u})
\]
if $u \in \domain{T_0}$ and $\abs{ \supp u} \geq r$.
Assume further that $\domain{T_0^K} = \domain{T_1^k}$
for $k \in \N$, $k \leq p$, $p > m/r\,+1$.
Then the wave operators $W_{\pm}(T_1,T_0)$ exist and
$\range(W_+) = \range(W_-) = \hilH_{ac,i}$.
\end{theorem}

% xxx proof of theorem 5.3

% xxx 5.4. Example
% xxx 5.5. Example
% xxx 5.6. Remark

\paragraph*{}
There are many related results concerning the examples 5.4 and 5.5;
some of thre are, at least partially, much stronger
(see for example [2,4,8,13,16,19]).
But it should be noted that Theorem 5.3 also covers perturbations which are
not differential operators. These results seem to indicate that strong
\textit{scattering results are possible with such smallness conditions near infinity}.

\paragraph*{}
Much weaker conditions are possible if one wants only to prove the existence
of wave operators (e.\,g. [15,24,25]).
There are even sufficient conditions [15,25] which do not require the relative
smallness at infinity. A very old result of this kind is due to Kuroda [15]
and says that for $T_0 = -\Delta$ and $T_1 = T_0 +q$ with
\begin{equation}\tag{5.7}
xxx
\end{equation}
the wave operators $W_{pm}(T_1,T_0)$ exist.
This result has been extended to more general operators [25].

\paragraph*{}
For $m \geq 3$ the multiplication by such a $q$ is (i.\,g.) not $T_0$-small
at infinity in general.
For an \textit{example} chose $q(x) \eq -1$ in the union of infinitely many
balls $B_n$ of radii $r_n$ with centers $x_n \in \R^m$.
If the sequence $(x_n)$ tends to infinity fast enough (depending on $(r_n)$),
then Kuroda's condition (5.7) is satisfied.
But if the sequence $(r_n)$ does not converge to $0$, then this operator
is not $T_0$-small at infinity: let $r_n \geq \rho_0$,
$u \in C_0^{\infty}(x \in \R^m: \abs{x} < \rho_0)$, $u \neq 0$,
$u_n(x) = u(x-x_n)$, then $q\,u_n = -u_n \neq 0$.

\paragraph*{}
Actually such perturbations do not even preserve the essential spectrum.
In order to see this we chose $\rho_0 > 0$ so large that there is a
$u \in C_0^{\infty}(\{x \in \R^m: \abs{x} < \rho_0\})$ such that
\[
xxx.
\]
Then we have for every $n$ (with the above notation)
\[
xxx
\]
and, since the $u_n$ have mutually disjoint support,
the same inequality holds for every linear combination.
Therefore there is an infinite-dimensional subspace $M$ such that this inequality
holds for every $u \in M$.
By theorem 3.6 we then have $\dim E(-\frac{1}{2}) = \infty$.
Since $-\Delta +q \geq -1$ holds, there must be at least one point of
$\sigma_e(T)$ in $[-1,-\frac{1}{2}]$.

\paragraph*{}
It is also easy to show that such conditions do not in general imply
$R(W_+) = R(W_-) = H_{ac,1}$.
We cannot prove this for $m \leq 3$.
For $m \geq 4$ we give the following counter example. Chose
\[
xxx.
\]
This function obviously satisfies Kuroda's condition (5.7).
But by means of a separation of variables it is easy to show
that, if $\rho$ is sufficiently large,
$\sigma(T_0+q) = [\mu,\infty)$ with some $\mu < 0$;
furthermore this operator is purely absolutely continuous.
Therefore the relation $R(W_+) = R(W_-) = H_{ac,1}$ cannot hold.

\paragraph*{}
This indicates that such weak conditions are not sufficient for a reasonable
scattering theory. The best possible result which one might expect is
the following:
"Let $T_0$ be a self-adjoint differential operator with constant coefficients.
Assume (i) $V$ is $T_0$-small at infinity, (ii) $T_1 = T_0+V$ is self-adjoint
and (iii) the wave operators $W_{\pm}(T_1,T_0)$ exist.
Then $R(W_+) = R(W_-) = H_{ac,1}$".

%------------------------------------------------------------------------------%
\section{Appendix (added after the end of the symposium)}                      %
%------------------------------------------------------------------------------%
In the course of this symposium the author was able to conclude from part ii)
of the proof of Theorem 3.4 that every singular sequence of $A+B$ and $\lambda$
is also a singular sequence of $A$ and $\lambda$;
together with part i) this means that under the assumption of Theorem 3.4
the operators $A$ and $A+B$ have the same singular sequences.
This simplifies several later results.

\subparagraph*{}
The symmetry in this result lead the author to the conjecture
that "$B$ is also $(A+B)^2$-compact" or more general:
"If $A_1$ and $A_2$ are self-adjoint with the same domain $\cal{D}$ and
$\cal{D} \subset \domain{B}$,
then $B$ is $A_1^2$-compact if and only if $B$ is $A_2^2$-compact".
The proof of the first conjecture is contained in the following theorem
which the authors owes to T.\,Kato.
So far we could neither prove nor disprove the second conjecture.
It is (in general) not true if the operators $A_j$ are not self-adjoint.

\begin{theo}
\textit{
(T.\,Kato, private communication).
Let $A_1$ and $A_2$ be self-adjoint with $\domain{A_1} = \domain{A_2}$.
Assume that $B: A_2 - A_1$ is $A_1^2$-compact. Then
\begin{itemize}
\item[a)]
$B$ is also $A_2^2$-compact.
\item[b)]
$\sigma_e(A_1) = \sigma_e(A_2)$
\item[c)]
an operator $V$ is $A_1$-bounded and $A_1^2$-compact \iff
$V$ is $A_2$-bounded and $A_2^2$-compact.
\end{itemize}
}
\end{theo}

\Proof:
Let $R_j = (i-A_j)^{-1}$ for $j=1,2$.
Then $B R_j$ are bounded operators in $\hilH$ for $j=1,2$ and $B R_1^2$
is a compact operator in $\hilH$.
We also have
\begin{equation}\tag{1}
R_2 - R_1 \= R_1 B R_2 - R_2 B R_1
\end{equation}
and therefore
\begin{equation}\tag{2}
R_2^2 \= (I + R_2 B) \, R_1^2 (I + B R_2).
\end{equation}

\paragraph{}
a) From (2) we get $BR_2^2 \eq (I + B R_2) \, B \, R_1^2 \, (I+BR_2).$
Hence $BR_2^2$ is compact because $BR_1^2$ is compact.
Hence $B$ is $A_2^2$-compact.

b) From (2) we also get
\[
R_2^2 \= R_1^2 -(BR_1^{*2})^*R_2 + R_2BR_1^2(I+BR_2)
\]
which is compact since $BR_1^{*2}$ and $BR_1^2$ are compact.
Hence $\sigma_e(R_2^2) = \sigma_e(R_1^2)$,
and therefore $\sigma_e(A_2) = \sigma(A_1)$.
(Here a simple spectral mapping theorem for the essential spectrum is used).

c) Let $V$ be $A_1$-bounded and $A_1^2$-compact.
Then by (1)
\[
VR_2 \= VR_1(I-BR_2)
\]
is a bounded operator in $\hilH$.
Therefore $V$ is $B$-bounded. By (2)
\[
VR_2^2 \= (VR_1^2 + VR_2BR_1^2)(I + BR_2)
\]
is compact since $VR_1^2$ nad $BR_1^2$ are compact.
Hence $V$ is $B^2$-compact.
The converse is proved in the same way since the assumptions are completely
symmetric {\wrt} $A_1$ and $A_2$ by part a) of this theorem. \qed

\paragraph{}
Part b) of the above theorem gives a very simple proof of this invariance
of the essential spectrum.
If we combine part a) of this theorem with part i) of the proof of theorem 3.4,
then we get a simple proof of the fact
that also the singular sequence are preserved.
The complicated part ii) of the proof of theorem 3.4. is not needed.

\paragraph{}
Note that part a) of the theorem holds also
if $A_j$ are arbitrary closed operators in a Banach space and
$\rho(A_1) \cap \rho(A_2) \neq \emptyset$.

%------------------------------------------------------------------------------%
%                                 BIBLIOGRAPHY                                 %
%------------------------------------------------------------------------------%
\newpage
\fancyhead[C] {}
\fancyhead[OL]{}
\fancyhead[ER]{}
\thispagestyle{empty}

\titleformat{\section}
{\normalfont \bfseries \filcenter}
{\thesection.\;}{0mm}{}

\newlength{\bibitemsep}\setlength{\bibitemsep}{.5\baselineskip plus .05\baselineskip minus .05\baselineskip}
\newlength{\bibparskip}\setlength{\bibparskip}{0pt}
\let\oldthebibliography\thebibliography
\renewcommand\thebibliography[1]{%
  \oldthebibliography{#1}%
  \setlength{\parskip}{\bibitemsep}%
  \setlength{\itemsep}{\bibparskip}%
}

\vspace{8mm}
\begin{thebibliography}{99}
%\setlength\bibitemsep{3\itemsep}
\begin{small}

\bibitem{1}  Aronszajn,\,N.:
On a problem of Weyl in the theory of singular Sturm Liouville equations.
Amer.\,J.\,Math. 79, 597-610 (1951).

\bibitem{2}  Birman,\,M.\,Sh.:
Scattering problems for differential operators with constant coefficients.
Functional Anal. Appl. 3, 167-180 (1969).

\bibitem{3}  B{\"o}cker,\,U.:
Invarianz des wesentlichen Spektrums bei {\SCHROEDINGER}-Operatoren.
Math. Ann. 207, 349-352 (1974).

\bibitem{4}  Eckardt,\,K.-J.:
Scattering Theory for Dirac Operators.
Math. Z. (to appear)

\bibitem{5}  Glazman,\,I.\,M.:
Direct Methods of Qualitative Spectral Analysis
of Singular Differential Operators,
Israel Program for Scientific Translations, Jerusalem 1965.

\bibitem{6}  Gustafson,\,K. and J.\,Weidmann:
On the essential spectrum.
J.\,Math.\,Anal.\,Appl. 25, 121-127 (1969).

\bibitem{7}  Ikebe,\,T. and T.\,Kato:
Uniqueness of the self-adjoint extensions
of singular elliptic differential operators.
Arch.\,Rational\,Mech.\,Anal. 2, 77-92 (1969).

\bibitem{8}  Ikebe,\,T. and T.\,Tayoshi:
Wave and Scattering Operators for Second Order Elliptic Operators in $\R^3$.
Publ. Res. Inst. Math. Sci.,
Kyoto University, Ser.A, 4, No.2, 483-496 (1968).

\bibitem{9}  {\JOERGENS},\,K.:
Wesentliche Selbstadjungiertheit singul{\"a}rer elliptischer
Differentialoperatoren zweiter Ordnung in C$_0^{\infty}$(G).
Math Scand. 15, 5-17 (1964).

\bibitem{10} {\JOERGENS},\,K.:
Perturbations of the Dirac Operator.
Conf. Ordinary and Partial Differential Equations, Dundee/Scotland 1972,
Lecture Notes in Math. 280, 87-102 (1972).

\bibitem{11} {\JOERGENS},\,K. and J.\,Weidmann:
Spectral Prperties of Hamitonian Operators.
Lecture Notes in Math., 313, 1973.

\bibitem{12}  Kato,\,T.:
Perturbation theory for linear operators.
Berlin-Heidelberg-New York:
Springer, 1966.

\bibitem{13}  -----:
Some results on potential scattering.
Proceedings at the International Conference on Functional Analysis
and Related Topics,
206-215, Tokyo, University of Tokyo Press, 1969.

\bibitem{14}  -----:
{\SCHROEDINGER} operators with singular potentials.
Israel J. Math. 13, 135-148 (1972).

\bibitem{15}  xxx,\,:

\bibitem{16}  xxx,\,:

\bibitem{17}  xxx,\,:

\bibitem{18}  xxx,\,:

\bibitem{19}  xxx,\,:

\bibitem{20}  xxx,\,:

\bibitem{21}  xxx,\,:

\bibitem{22}  xxx,\,:

\bibitem{23}  xxx,\,:

\bibitem{24}  xxx,\,:

\bibitem{25}  xxx,\,:

\bibitem{26}  xxx,\,:

\bibitem{27}  Weidmann,\,J.:
Zur Spektraltheorie von Sturm-Liouville-Operatoren.
Math.\,Z. 28, 268-302 (1967).

\bibitem{28}  Weidmann,\,J.:
Perturbations of self-adjoint operators in $L_2(G)$ with applications
to differential operators.
Conf. Ordinary and Partial Diff. Equations Dundee/Scotland 1974 (to appear in 
Lecture Notes in Math.)

\bibitem{29}  Yosida,\,K.:
Functional Analysis.
Berlin-Heidelberg-New York:
Springer, 1966.

\end{small}
\end{thebibliography}
%------------------------------------------------------------------------------%
%                               END OF DOCUMENT                                %
%------------------------------------------------------------------------------%
\end{document}
